%----------------------------------------------------------------------
\section{Homotopical and Categorical Structure}\label{sec:homotopy}
%----------------------------------------------------------------------

The grading of the exterior algebra provides a direct\footnote{Section-level Toulmin. \emph{Move}: the grade structure of $\bigwedge V$ yields chain complexes and groupoids. \emph{ZFC comparison}: ZFC has no native homotopical structure; it must be imported via simplicial sets or CW-complexes. Here, $\partial \circ \partial = 0$ is a theorem of $\GF(2)$ arithmetic. \emph{Residue instance}: homological gaps ($H_k \neq 0$) are the homotopy-level manifestation of the same residue mechanism that produces gaps/overlaps at the set-theoretic level. The mechanism is uniform: residue = $\chi = 0$ at every grade.}
correspondence with homotopy theory. The key observation
is that the framework has exactly one primitive --- the
wedge product over $\GF(2)$ --- and all categorical
structure (objects, morphisms, higher morphisms, identity,
composition, boundaries) is this single operation applied
to different inputs.

\subsection{The single-primitive principle}

\begin{definition}[Degeneracy as identity]\label{def:degeneracy}
A \emph{0-cell} is not a primitive distinct from morphisms.
A 0-cell is a 1-cell (discriminator) whose source and target
coincide --- that is, a discriminator applied to itself:
\begin{equation}\label{eq:object-as-loop}
  d \wedge d = 0.
\end{equation}
The ``object'' associated with discriminator $d$ is the
degenerate self-discrimination of $d$. Grade 0 is not
``scalars, interpreted as objects.'' Grade 0 is \emph{what
results when a discrimination collapses}.
\end{definition}

\begin{remark}
This subsumes the last remaining ontological primitive.
The framework does not say ``there are objects and there are
morphisms.'' It says: \emph{there is only discrimination.}
The grade of an element measures its degree of
non-degeneracy. Grade 0 is fully degenerate
(self-referential, annihilated). Grade 1 is a single
non-degenerate discrimination. Grade $k$ is $k$ independent
non-degenerate discriminations composed. Objects are the
fixed points of the self-discrimination operation.
\end{remark}

\begin{definition}[Discrimination complex]\label{def:complex}
The \emph{discrimination complex} of a regime $\kappa$ with
$\dim(\kappa) = n$ is the exterior algebra
$\bigwedge(\GF(2)^n) = \bigoplus_{k=0}^{n} \bigwedge^k(\GF(2)^n)$,
with the grading interpreted as:
\begin{itemize}[nosep]
  \item \textbf{Grade 0}: degenerate (self-)discriminations;
    equivalently, $\kappa$-equivalence classes of elements.
  \item \textbf{Grade $k$} (for $1 \leq k \leq n$):
    $k$-fold non-degenerate compound discriminations, i.e.\
    $k$-morphisms.
\end{itemize}
The complex is truncated at grade $n$, since
$\bigwedge^k(\GF(2)^n) = 0$ for $k > n$. An $n$-dimensional
$\kappa$ yields an $n$-groupoid.
\end{definition}

\subsection{Boundary maps}\label{sec:boundary}

The identification of 0-cells with degenerate 1-cells gives
a natural definition of boundary: the boundary of a $k$-cell
is obtained by contracting one factor, reducing the grade by
one.

\begin{definition}[Boundary operator]\label{def:boundary}
The \emph{boundary operator}
$\partial_k: \bigwedge^k V \to \bigwedge^{k-1} V$ is the
linear map defined on basis elements by
\begin{equation}\label{eq:boundary}
  \partial(e_{i_1} \wedge \cdots \wedge e_{i_k})
  = \sum_{j=1}^{k}
    e_{i_1} \wedge \cdots \wedge \widehat{e_{i_j}} \wedge
    \cdots \wedge e_{i_k}
\end{equation}
where $\widehat{e_{i_j}}$ denotes omission and the sum
is in $\bigwedge^{k-1} V$ (a formal sum over $\GF(2)$,
not a set). The map is extended to all of $\bigwedge^k V$
by $\GF(2)$-linearity. For a 1-cell $e_i$,
$\partial(e_i) = 1 \in \bigwedge^0 V = \GF(2)$ (the
unit). For a 0-cell, $\partial = 0$.
\end{definition}

\begin{proposition}[Chain complex]\label{prop:boundary}
The boundary operator satisfies:
\begin{enumerate}[label=(\roman*),nosep]
  \item $\partial_k$ is a $\GF(2)$-linear map
    $\bigwedge^k V \to \bigwedge^{k-1} V$.
  \item The boundary of a 2-cell $e_1 \wedge e_2$ is
    $e_2 + e_1 = e_1 + e_2$: the formal sum (over $\GF(2)$)
    of its two constituent 1-cells.
  \item $\partial_{k-1} \circ \partial_k = 0$ for all $k$.
    The graded sequence
    $\cdots \xrightarrow{\partial_3} \bigwedge^2 V
    \xrightarrow{\partial_2} \bigwedge^1 V
    \xrightarrow{\partial_1} \GF(2) \to 0$
    is a chain complex over $\GF(2)$.
\end{enumerate}
\end{proposition}

\begin{proof}
(i) is by construction: $\partial$ is defined on basis
elements and extended by linearity.

(ii): $\partial(e_1 \wedge e_2) = e_2 + e_1$ by
equation (\ref{eq:boundary}).

(iii): Apply $\partial$ twice to a basis $k$-cell
$\omega = e_{i_1} \wedge \cdots \wedge e_{i_k}$:
\[
  \partial(\partial(\omega))
  = \partial\!\left(\sum_{j=1}^{k}
    e_{i_1} \wedge \cdots \wedge
    \widehat{e_{i_j}} \wedge \cdots \wedge e_{i_k}\right)
  = \sum_{j=1}^{k} \sum_{\ell \neq j}
    e_{i_1} \wedge \cdots \wedge \widehat{e_{i_j}}
    \wedge \cdots \wedge \widehat{e_{i_\ell}}
    \wedge \cdots \wedge e_{i_k}.
\]
Each $(k{-}2)$-cell (obtained by omitting $e_{i_j}$ and
$e_{i_\ell}$) appears exactly twice in this double sum:
once from omitting $j$ first then $\ell$, once from
omitting $\ell$ first then $j$. Over $\GF(2)$,
$1 + 1 = 0$, so every term cancels.
$\partial \circ \partial = 0$.
\end{proof}

\begin{remark}
The identity $\partial \circ \partial = 0$ is the defining
property of a chain complex. The discrimination complex is
therefore a \emph{chain complex over $\GF(2)$}, and its
homology groups $H_k = \ker \partial_k / \operatorname{im}
\partial_{k+1}$ are well-defined. These homology groups
classify the ``holes'' in the discrimination regime --- the
residues that have not yet been consumed. This connects the
algebraic residue concept (Definition \ref{def:residue})
to the topological notion of homological holes.
\end{remark}

\begin{theorem}[Exhaustivity is operational]\label{thm:exhaustive}\footnote{Axiom class \textsf{A} $\cdot$ depth 3 $\cdot$ \S\ref{app:exhaustivity-full}.  \emph{Warrant}: $\partial \circ \partial = 0$ (double-counting over $\GF(2)$). The residue is fully characterized at detection. Finite-$\kappa$: yes. Novel (residue = homological gap, operationalized).}
The identity $\partial \circ \partial = 0$ entails that every
residue fully specifies its own shape. That is: the homology
class of a residue constitutes a complete description of the
discriminator required to consume it.
\end{theorem}

\begin{proof}
A residue $R_d(S)$ (Definition \ref{def:residue}) consists of
elements in $\ker \chi_d$ --- those invisible to discriminator
$d$. In the chain complex, these elements occupy a cycle in
$\ker \partial_k$ that is not in $\operatorname{im} \partial_{k+1}$:
they are detected as a non-trivial homology class $[R] \in H_k$.

The identity $\partial \circ \partial = 0$ means that the
boundary of a boundary is empty. In discrimination terms:
\emph{the residue of a residue is always zero.} If the
discrimination process has identified a gap, there is no
further gap-ness to discover about that identification. The
gap is fully characterized at the moment of detection.

Concretely, the homology class $[R]$ specifies:
\begin{enumerate}[label=(\roman*),nosep]
  \item the grade $k$ at which the gap exists (the level of
    the homotopy tower where discrimination is incomplete),
  \item the specific $(k-1)$-cells on the boundary of the
    missing $k$-cell (the existing discriminations that
    bound the gap), and
  \item the precise shape of the $k$-cell required to fill
    the gap, if such a cell exists. When $H_k(\kappa) = 0$
    under the full regime (all cycles are boundaries),
    the filling is unique. When $H_k(\kappa) \neq 0$,
    the residue specifies a homology class but not a
    unique filling. The claim that the filling is unique
    under the full $\kappa$ is equivalent to the claim
    that $H_k(\kappa) = 0$ for the relevant $k$, which
    is a structural property of the specific regime,
    not a general theorem.
\end{enumerate}

Therefore $\kappa$-evolution is not a search process. Each
residue arrives with the specification of the discriminator
needed to consume it. The homology class \emph{is} the
blueprint for the cleavage. (Full argument: \S\ref{app:exhaustivity-full}.)
\end{proof}

\begin{corollary}[Non-uniqueness under truncation]
\label{cor:truncation-ambiguity}
Under a truncated (projected) $\kappa$-regime $\kappa'$,
a cycle may admit multiple distinct fillings. The
multiplicity of fillings is not a diagnostic \emph{about}
a residue --- it \emph{is} the residue. Specifically, the
wedge product of the candidate filling cells under the
truncated regime produces a non-zero element: this element
is the residue of the truncation, and its detection is
itself a $\kappa$-evolution trigger.
\end{corollary}

\begin{proof}
Let $\kappa' \subset \kappa$ be a sub-regime obtained by
omitting dimensions. The chain complex of $\kappa'$ is a
quotient of the chain complex of $\kappa$: some boundary
relations are forgotten. A cycle $[R]$ that admits a unique
filling in $H_k(\kappa)$ may, upon projection, map to a
class in $H_k(\kappa')$ that is compatible with multiple
fillings $c_1, \ldots, c_m$, since the forgotten dimensions
carried the constraints that distinguished them.

The candidate fillings $c_1, \ldots, c_m$ are distinct cells
that agree on their boundary in $\kappa'$ but differ in
the omitted dimensions. The wedge product $c_i \wedge c_j$
for $i \neq j$ is therefore non-zero in the full exterior
algebra: it detects a discrimination that exists in $\kappa$
but is invisible in $\kappa'$. This non-zero wedge product
is a residue (Definition \ref{def:residue}) --- an element
in the kernel of the XOR signature under the truncated
regime.

Crucially, computing the wedge product of the candidate
fillings to detect their non-uniqueness \emph{is itself a
discrimination}. The act of observing the ambiguity is a
wedge product that produces the residue that triggers the
$\kappa$-evolution needed to resolve it. The observer is
not independent of the system: detecting that a projection
is too coarse is the same act as beginning to refine it.
\end{proof}

\begin{remark}
This is a stronger claim than ``homology detects holes.''
The standard topological statement is descriptive: homology
classifies gaps passively. The discrimination framework
makes exhaustivity \emph{operational}: because every
discrimination either cleanly partitions (producing no
residue) or produces a residue whose shape is fully
determined by $\partial \circ \partial = 0$, the system
always knows the exact structure of what it cannot yet
discriminate. Gaps are not merely detected but
\emph{self-specifying}.
\end{remark}

\subsection{Groupoid structure from $\GF(2)$}

\begin{theorem}[Intrinsic groupoid structure]\label{thm:groupoid}\footnote{Axiom class \textsf{A} $\cdot$ depth 3 $\cdot$ \S\ref{app:groupoid-full}.  \emph{Warrant}: commutativity of $\wedge$ over $\GF(2)$ ($-1 = 1$). Every 2-cell is its own inverse. Finite-$\kappa$: yes. Standard (chain complexes yield groupoids). Model: $\wedge$ verified in \S\ref{sec:concrete-model}.}
The discrimination complex $\bigwedge(\GF(2)^n)$ is a
symmetric $n$-groupoid in which every $k$-morphism is its own
inverse, objects are degenerate morphisms, and boundaries are
computed by contraction.
\end{theorem}

\begin{proof}
We verify the groupoid axioms, noting that all structure
derives from the wedge product.

\emph{Objects.} 0-cells are elements of $\bigwedge^0 V \cong
\GF(2)$. By Definition \ref{def:degeneracy}, these are
degenerate 1-cells: $d \wedge d = 0$ for any discriminator $d$.
Objects are not a separate primitive.

\emph{Self-inverse property.} Over $\GF(2)$, every element $x$
satisfies $-x = x$, since $\operatorname{char}(\GF(2)) = 2$.
Therefore every $k$-morphism is its own additive inverse.
Traversing a discrimination and reversing it are the same
operation.

\emph{Contraction.} For any grade-1 element $d$,
$d \wedge d = 0$ (equation \ref{eq:self-wedge}). A path
composed with itself is contractible --- homotopically trivial.
This extends to higher grades: for any $\omega \in \bigwedge^k V$,
the component of $\omega \wedge \omega$ in $\bigwedge^{2k} V$
vanishes when expanded in a basis, since every term contains
a repeated factor.

\emph{Composition.} The wedge product provides composition at
every level. For grade-1 elements $d_1, d_2$, the composite
$d_1 \wedge d_2 \in \bigwedge^2 V$ is the 2-morphism
connecting $d_1$ and $d_2$ (Proposition \ref{prop:boundary},
claim (ii)). Associativity is inherited from the exterior
algebra:
$(d_1 \wedge d_2) \wedge d_3 = d_1 \wedge (d_2 \wedge d_3)$.

\emph{Boundaries.} The contraction boundary
(Definition \ref{def:boundary}) provides source and target maps,
with $\partial \circ \partial = 0$
(Proposition \ref{prop:boundary}).

\emph{Symmetry.} Over $\GF(2)$, for $\alpha \in \bigwedge^j V$
and $\beta \in \bigwedge^k V$,
$\alpha \wedge \beta = (-1)^{jk} \beta \wedge \alpha
= \beta \wedge \alpha$ since $-1 = 1$ in $\GF(2)$. The
groupoid is symmetric at every level.

These properties --- degenerate objects, self-inverse
morphisms, contractible self-composition, associative and
symmetric composition, and well-defined boundaries ---
constitute a symmetric $n$-groupoid $\bigwedge(\GF(2)^n)$. (Axiom-by-axiom check: \S\ref{app:groupoid-full}.) Derived entirely from
$\bigwedge(\GF(2)^n)$.
\end{proof}

\begin{remark}
The symmetry deserves emphasis. The order of discrimination
does not matter at any level of the homotopy tower. This is
the homotopical expression of the order-independence result
(Proposition \ref{prop:monoidal}).
\end{remark}

\subsection{Relationship to homotopy type theory}

The discrimination complex admits a natural reading in the
language of homotopy type theory (HoTT).

\begin{proposition}[Univalence as extensionality]\label{prop:univalence}
The univalence axiom of HoTT --- that equivalent types are
equal --- corresponds to Theorem \ref{thm:ext}
($\kappa$-extensionality). Two $\kappa$-equivalence classes
that cannot be distinguished by any discriminator in $\kappa$
are \emph{equal as types}. The equivalence is not merely
logical but homotopical: the path space between them (the
set of discriminators that could separate them) is
contractible (empty, hence trivially contractible).
\end{proposition}

\begin{remark}[The operationalist axiom as native univalence]
\label{rem:native-univalence}
In HoTT, univalence is an axiom: the canonical map from
identity to equivalence is an equivalence. In the
discrimination framework, the operationalist axiom
(Definition \ref{def:operationalist}) plays this role:
$\kappa$-indistinguishable $=$ identical. The structural
content is the same --- operational equivalence implies
identity --- but the mechanism differs. HoTT requires
constructing an equivalence (a map with contractible
fibers). The framework requires the \emph{absence} of a
separating discriminator, which is a negative condition.
The path space is contractible not because a contraction
was constructed but because there is no structure to hold
it open (Remark \ref{rem:hyper-recursive}).

The precise statement (Appendix B.36): the path space
between $a$ and $b$ is $\{a + b\} \cap \kappa^\perp$.
The operationalist axiom says $\bigcup \kappa = V$,
so $\bigcap \kappa^\perp = V^\perp = \{0\}$.
Persistent equivalence forces $a + b = 0$, hence
$a = b$. The ua map (equivalence $\to$ identity) and
idtoeqv (identity $\to$ equivalence) are inverses.
The path space is contractible. Univalence is
Galois-group exhaustion: the only persistent orbit
is the trivial orbit.
\end{remark}

\begin{remark}[Write path as type construction, read path
  as type elimination]\label{rem:hott-paths}
The write/read distinction (Section \ref{sec:write-read})
maps onto the introduction/elimination distinction in type
theory:
\begin{itemize}[nosep]
  \item The \emph{write path} (linear, residue-consuming)
    is \emph{type introduction}: constructing new types by
    attaching cells to the discrimination complex.
  \item The \emph{read path} (monoidal, freely iterable)
    is \emph{type elimination}: computing with existing
    types by querying the $\tau/\sigma$ structure.
\end{itemize}
That the write path is linear (each residue consumed once)
while the read path admits the exponential modality
(unlimited re-use) reflects the standard linear-logic
decomposition underlying type-theoretic computation.
\end{remark}

\begin{remark}[Transport as reclassification]
In HoTT, transport along a path moves inhabitants of one
type to another. In the discrimination framework, transport
along a grade-1 morphism (discriminator $d$) is the
reclassification of elements: an element classified under
$\tau_d$ is ``transported'' to the $\tau$-side of the
partition; an element classified under $\sigma_d$ to the
$\sigma$-side. The four-cell structure
(Definition \ref{def:discriminator}) generalizes transport
by allowing partial transport (gap) and over-determined
transport (overlap).
\end{remark}

\subsection{The constructed groupoid}

Combining the set-theoretic and homotopical structure, we
can now state the categorical content of the framework
more precisely than a 1-categorical account would allow.

\begin{definition}[Emergent $n$-groupoid]\label{def:n-groupoid}
Under a fixed regime $\kappa$ of dimension $n$, the
\emph{$n$-groupoid} $\mathcal{G}_\kappa$ is the
symmetric $n$-groupoid whose:
\begin{itemize}[nosep]
  \item 0-cells are $\kappa$-equivalence classes,
  \item $k$-cells (for $1 \leq k \leq n$) are non-zero
    elements of $\bigwedge^k(\GF(2)^n)$,
  \item composition at all levels is the wedge product,
  \item identity at all levels is witnessed by
    $\omega \wedge \omega = 0$.
\end{itemize}
This structure is constructed, not postulated: it is
a groupoid compatible with $\bigwedge(\GF(2)^n)$.
(Uniqueness of the symmetric groupoid: all groupoid
over the same algebra may exist.)
\end{definition}

\begin{proposition}[Growth of homotopical complexity]\label{prop:growth}
If $\kappa \to \kappa'$ by evolution (Definition
\ref{def:evolution}), then $\mathcal{G}_{\kappa'}$ is an
$(n+1)$-groupoid extending $\mathcal{G}_\kappa$. The
inclusion $\mathcal{G}_\kappa \hookrightarrow
\mathcal{G}_{\kappa'}$ preserves all existing cells and
adds new cells at every grade $k \leq n+1$ (those involving
the new discriminator $d_{n+1}$).
\end{proposition}

\begin{proof}
$\kappa'$ has dimension $n+1$, so
$\bigwedge(\GF(2)^{n+1})$ has non-trivial components through
grade $n+1$. The natural inclusion
$\GF(2)^n \hookrightarrow \GF(2)^{n+1}$ (appending a zero
in the $(n+1)$-th coordinate) induces an inclusion of
exterior algebras
$\bigwedge(\GF(2)^n) \hookrightarrow \bigwedge(\GF(2)^{n+1})$
that preserves the wedge product and hence all groupoid
structure. The new cells are precisely those involving the
basis vector $e_{n+1}$ corresponding to discriminator
$d_{n+1}$.
\end{proof}

\begin{openquestion}[Adjunctions from $\kappa$-deformation]
A $\kappa$-deformation that projects from a finer regime to
a coarser one induces a functor
$\mathcal{G}_\kappa \to \mathcal{G}_{\kappa'}$ (a forgetful
functor, forgetting discrimination dimensions). A
$\kappa$-deformation that embeds a coarser regime into a
finer one induces a functor in the opposite direction (a
free functor, freely adding undetermined cells). Under what
conditions do these form an adjoint pair? If they do, the
free-forgetful adjunction --- the fundamental construction
of applied category theory --- arises from the direct
summand structure that discrimination over $\GF(2)$
provides.
\end{openquestion}

